\documentclass[11pt]{article}
\usepackage[a4paper,margin=25mm]{geometry}
\usepackage{amsmath,amssymb,amsthm}
\usepackage[hidelinks]{hyperref}
\newtheorem{theorem}{Theorem}
\begin{document}
\title{The four-cycle contradicts the minimality claim in Conjecture 3461}
\author{Chengzhe Gao}\date{}\maketitle

The conjecture includes the claim that the smallest planar graph for which
$\chi_{\rm DP}-\chi=1$ is a three-dimensional prism. The four-cycle already
has this difference and has only four vertices. A prism over a genuine
$m$-gon has $2m\geq6$ vertices, so the asserted minimality is false.

Let $C_4$ have successive vertices $0,1,2,3$ and edges $01,12,23,30$.
A DP cover assigns a color fiber to each vertex, makes each fiber a clique,
and places a matching between fibers of each graph edge. There are no
cross-fiber edges over a nonedge. A DP coloring is an independent transversal,
choosing one color from every fiber. These are the standard definitions;
see \href{https://arxiv.org/abs/1609.00763}{Bernshteyn, Kostochka, and Pron,
On DP-coloring of graphs and multigraphs}.

\begin{theorem}
$C_4$ is planar, $\chi(C_4)=2$, and $\chi_{\rm DP}(C_4)=3$.
\end{theorem}
\begin{proof}
The usual alternating coloring of the cycle uses two colors. An edge
precludes a one-coloring, so the ordinary chromatic number is two.

To rule out DP two-colorability, give every vertex the fiber $\{0,1\}$.
On edges $01,12,23$, match equal colors. On edge $30$, instead match
opposite colors. Avoiding the first three matchings forces successive
chosen colors to alternate, so vertices $0$ and $3$ receive opposite
colors. Their chosen pair is then forbidden by the last matching. Thus
this genuine two-fold cover has no independent transversal.

For the upper bound, take an arbitrary cover with at least three colors
available at each vertex, and restrict each fiber to any three of them.
Choose a color at vertex $0$. At vertex $1$ at most one color is forbidden
by the previous choice, so choose another. At vertex $2$ again at most
one is forbidden. At vertex $3$ its two already colored neighbors forbid
at most two colors in total. One of its three colors therefore remains.
This constructs an independent transversal for every cover, proving
DP three-colorability. Covers with empty fibers cannot have a transversal,
and a one-fold cover with its edge matchings present also has none; hence
the exact DP chromatic number is three.

Finally place the vertices at $(0,0),(1,0),(1,1),(0,1)$ in that order.
Draw each edge as the corresponding side of the unit square. Its sides
have disjoint interiors and meet only at the prescribed shared endpoints,
so this is a planar embedding.
\end{proof}

The difference is $3-2=1$ on four vertices. This is smaller than any prism
over a polygon; the four-cycle also has fewer edges. The counterexample
targets the explicit minimality clause, and makes no claim to disprove
the other planar DP-coloring bounds in the source. In particular no girth
five hypothesis is imposed on the source's separate difference-spectrum
and minimality clause.

\paragraph{Formal verification of covers and chromatic numbers.}
The Lean graph has the exact cyclic adjacency relation, with symmetry and
absence of loops checked. Its cover structure uses arbitrary vertex-dependent
color types and arbitrary symmetric, supported matchings; the matchings
may be partial. The auxiliary graph on the disjoint union of fibers is
defined with the required fiber cliques and cross-fiber edges. The program
proves that avoiding all forbidden pairs is equivalent to an independent
transversal in this graph.

The twisted two-fold cover satisfies every cover axiom. All $16$ possible
color selections are checked, and a proved tuple expansion shows that this
checks every function on the four vertices. The upper-bound proof is instead
symbolic in an arbitrary cover. A small three-color lemma proves that two
sets each containing at most one forbidden color cannot exclude all three
choices. Restriction from arbitrary color fibers to any three distinct
colors is explicitly formalized, with its injectivity and matching bridge.
Consequently the upper bound covers all lists of size at least three, not
only one selected cover or uniform named color lists. Both chromatic
parameters are proved minimal over every smaller natural number of colors.

\paragraph{Formal verification of the drawing and final contradiction.}
The square certificate defines the vertex positions, all four graph-edge
endpoint pairs, and parametrized edge maps
$(t,0),(1,t),(t,1),(0,t)$. Endpoint incidence, injectivity of vertices and
edge maps, exclusion of other vertices from edge interiors, and disjointness
of distinct edge interiors are all proved. The proofs quantify over
arbitrary scalar types, distinct endpoint coordinates, and arbitrary
interior parameters. In particular they apply to every real $t,u\in(0,1)$;
the proof is not an integer-coordinate sampling test. Over the reals these
explicit maps are the ordinary continuous straight segments, furnishing
the standard planar drawing.

Finally the program proves that $2m>4$ whenever $m\geq3$, and negates the
necessary instance of the proposed prism minimality when confronted with
this square drawing and the exact chromatic numbers. Lean 4.19.0 checks
the package using \texttt{Std}, with no admitted declarations, custom
axioms, or \texttt{native\_decide}.
\end{document}
